%! Author = lili
%! Date = 4/27/26

% Skript Kapitel 3


\chapter{\gls{zv}n und ihre \gls{verteilung}}\label{ch:zufallsvariablen-und-ihre-verteilung}
\untit{Vorlesung 09.05.2025 (V04 2025)}
\gls{zv}n ermöglichen es uns nicht nur, Ereignisse übersichtlicher darzustellen, sondern sind auch
das optimale Werkzeug, um in systematischer Weise von einem \gls{wraum} mit vielen
Elementarereignissen zu einem kleineren \gls{wraum} zu wechseln.


\section*{Lernziele}
\begin{itemize}
    \item Konzepte
    \begin{itemize}
        \item \gls{zv}n als Abbildung und ihre \gls{verteilung} als \gls{wmaß}
        \item \gls{verteilungsfunktion_zv}
        \item gemeinsame \gls{verteilung} von \gls{zv}n, Randverteilung, Randdichte
    \end{itemize}
    \item Fakten
    \begin{itemize}
        \item Eigenschaften von \gls{verteilungsfunktion_zv}en
        \item Die gemeinsame \gls{verteilung} bestimmt die Randverteilungen.
        \item Die Umkehrung ist im allgemeinen falsch.
    \end{itemize}
\end{itemize}


\section{Die \gls{verteilung} einer \gls{zv}n}\label{sec:die-verteilung-einer-zufallsvariablen}
\begin{bsp}[n-faches Werfen einer fairen Münze]
    \label{bsp:n-faches-Werfen-einer-fairen-Muenze} \txc{(2025)}  \\
    \[
        \glssymbol{eraum}_n = \kl{0,1}^n, \quad \glssymbol{wfk}_n (\kl{\glssymbol{elementarereignis}}) = \left(\frac{1}{2}\right)^n \faoio \txt{(Gleichverteilung)}
    \]
    $\glssymbol{elementarereignis}_i = 1$ bedeute ``Kopf'' \\
    Falls es uns interessiert, wie oft ``Kopf'' fällt:
    \[
        \sn(\glssymbol{elementarereignis}) \coloneqq  \sum^n_{i=1}     \glssymbol{elementarereignis}_i \txt{ist Anzahl ''Erfolge'' (Kopf)}
    \]
    \[
        \sn :\glssymbol{eraum} _n \rightarrow \kl{0,1,..,n} \txt{Abbildung!}
    \]
    Falls der relative Anteil der Erfolge relevant ist:
    \[
        T_n: \glssymbol{eraum} _n \rightarrow [0,1]
    \]
    \[
        T_n (\glssymbol{elementarereignis}) \coloneqq  \frac{1}{n} \sn (\glssymbol{elementarereignis}) = \frac{1}{n} \sum^n_{i=n} \glssymbol{elementarereignis}_i \faoio
    \]
\end{bsp}

%%%%%%%%%%%%%%%%%%%%%
% ZV %
%%%%%%%%%%%%%%%%%%%%%

\begin{defi}[\gls{zv}]
    \label{def:zv}\txc{(2025)}  \\
    Für einen \gls{wraum} $\wraum$ und eine Menge $E \neq \emptyset$ nennen wir eine messbare Abbildung
    \[
        \zvx: \glssymbol{eraum}  \rightarrow E
    \]
    eine $E$-wertige \gls{zv}.
\end{defi}

% TODO definition e-wertig und erklärung e-wertig

\komm{\defirefman{zv}: \gls{zv} ist keine Variable: Es ist eine Zuordnungsvorschrift der Stochastik, welche jedem möglichen Ergebnis eines Zufallsexperiments eine Größe zuordnet. Also das Ergebnis eines Zufallsexperiments.}

\begin{bem}
    \txc{(Skript 26 - Bemerkung 3.3.)}  \\
    Ist der Ereignisraum $\eraum$ abzählbar, so ist eine \gls{zv} nichts anderes als eine Abbildung, die auf
    dem $\eraum$ des \gls{wraum}s definiert ist.
    Ist $\eraum \subseteq \rz_d$, kommen eigentlich noch
    Messbarkeitsfragen hinzu, die wir im Rahmen dieser Vorlesung jedoch ignorieren wollen.

    Wie wir bereits in Beispiel \autoref{bsp:n-faches-Werfen-einer-fairen-Muenze} gesehen haben, ist es praktisch,
    Ereignisse mit Hilfe von \gls{zv}n auszudrücken.
    Wenn uns nur Ereignisse interessieren, die durch ein und dieselbe \gls{zv}
    ausgedrückt werden können, so können wir sogar noch einen Schritt weitergehen und einen neuen
    \gls{wraum} einführen, der nur noch solche Ereignisse kennt.
    Die folgende Definition ist
    ein erster Schritt in diese Richtung.
\end{bem}

\begin{defi}[Bild- und Wertebereich einer \gls{zv}n] \txc{(Skript 26 - Definition 3.4.)}  \\
    Sei $\zvx : \eraum \rightarrow E$ eine E-wertige \gls{zv} auf einem \gls{wraum} $\wraum$.
    Dann
    heißt E der Wertebereich von $\zvx$, und
    \[
        \zvx(\eraum) = {\zvx(\ele) : \ele \in \eraum}
    \]
    der Bildbereich von $\zvx$.
\end{defi}

\begin{bem}
    \txc{(Skript 26 - Bemerkung 3.5.)}  \\
    Es gilt stets $\zvx(\eraum) \subseteq E$.
    Der Bildbereich ist die kleinstmögliche Menge, die alle Werte enthält,
    die $\zvx$ annehmen kann.

    st Sn wie in Beispiel \autoref{bsp:n-faches-Werfen-einer-fairen-Muenze} die Anzahl der Erfolge in einem
    Zufallsexperiment der Länge $n \in \nz$,
    so können wir Sn auffassen als Abbildung von $\eraum$ nach $E = \nz_0$, aber auch als Abbildung
    nach $E = \{0, 1, \dots , n\}$ oder gar nach $E = \rz$.
    Unabhängig davon, wie wir den Wertebereich E
    wählen, gilt für den Bildbereich stets $\zvx(\eraum) = \{0, 1, \dots , n\}$.
\end{bem}

%%%%%%%%%%%%%%%%%%%%%
% \gls{verteilung} %
%%%%%%%%%%%%%%%%%%%%%

\begin{defi}[\gls{verteilung} einer \gls{zv}]
    \label{def:verteilung-zv} \txc{(2025)}  \\
    Sei $\zvx$ eine $E$-wertige \gls{zv} auf $\wraum$
    \[
        \pex(B) \coloneqq  \pe(\kl{\oio : \zvx(\glssymbol{elementarereignis}) \in B}) \quad \forall \txt{messbaren } B \subset E = \oub{\pe(\zvx \in B)}{Kurzschreibweise}
    \]
    Das \gls{wmaß} $\pex$ auf $E$ heißt \gls{verteilung} von $\zvx$.
\end{defi}
\komm{\defirefman{verteilung-zv}: Die \gls{verteilung} von $\zvx$ beschreibt, wie die Werte von $\zvx$ verteilt sind, also
\textbf{wie wahrscheinlich verschiedene Ergebnisse von $\zvx$ sind}.}

\begin{bsp}[Fortsetzung Münzwurf]
    \label{bsp:n-faches-Werfen-einer-fairen-Muenze-2}  \txc{(Skript 26 - Beispiel 3.7.)}  \\
    Die \gls{zv} $\sn$ aus Beispiel \autoref{bsp:n-faches-Werfen-einer-fairen-Muenze} hat den Bildbereich $\sn(\eraum_n)  = \{0, \dots , n\}$.
    Für $k \in \sn(\eraum_n)$ gilt
    \[
        \pe_{\sn}(\{k\}) = \pe_n (\{\ele \in \eraum_n : \sn (\ele) \in \{k\}\}) = \pe_n \left(  \left\{ \ele \in \eraum_n : \sumEN \ele_i = k \right\} \right) = \binom{n}{k}2^n
    \]
    weil $\pe_n$ die Gleichverteilung auf $\eraum_n = \{0,1\}^n$ ist und es $\binom{n}{k}$ viele Elementarereignisse $\ele \in \eraum_n$ gibt, für die $\sn (\ele) = k$ gilt.

    Ist $B \subseteq \sn(\eraum_n)$, so folgt mit den Rechenregeln für Wahrscheinlichkeiten, dass:
    \begin{align*}
        \pe_{\sn}(B) & = \pe_n(\{\ele \in \eraum_n : \sn (\ele) \in B\}) \\
        & = \pe_n \left(  \left\{ \ele \in \eraum_n : \sumEN \ele_i \in B \right\} \right) \\
        & = \pe_n \left( \uplus_{k \in B} \left\{ \ele \in \eraum_n : \sumEN \ele_i = k  \right\} \right) \\
        & = \sum_{k \in B} \pe_n \left( \left\{ \ele \in \eraum_n : \sumEN \ele_i = k  \right\} \right) \\
        & = \sum_{k \in B} \binom{n}{k}2^n
    \end{align*}
\end{bsp}

\begin{bem}[Sprunghöhen der \gls{verteilungsfunktion_zv}]
    \txc{(Skript 26 - Bemerkung 3.8.)}  \\
    In Definition \autoref{def:verteilung-zv} wird behauptet, dass die Abbildung $\pex$, die auf den messbaren
    Teilmengen von E definiert ist und Werte in $[0, 1]$ annimmt, tatsächlich ein \gls{wmaß} ist.
    Das müssen wir natürlich zeigen.
    Wir werden uns das Leben aber ein wenig leichter machen, und den Beweis nur in
    dem Fall führen, in dem $\pe$ durch eine Wahrscheinlichkeitsfunktion $\wfk$ gegeben ist.

    In diesem Fall entfällt auch die Notwendigkeit zu prüfen, ob eine Menge $B \subseteq E$ messbar ist, denn
    wenn $\eraum$ eine abzählbare Menge ist, dann kann die \gls{zv} $\zvx$, die auf $\eraum$ definiert ist, ebenfalls
    nur abzählbar viele Werte annehmen.
\end{bem}

\begin{lem}
    \txc{(2025)} \\
    Ist $\glssymbol{eraum} $ abzählbar und $\pe$ durch die \gls{wfk} $\glssymbol{wfk}$ gegeben, so ist $\pex$ (aus Definition \autoref{def:verteilung-zv})
    ein \gls{wmaß}.

    \paragraph{Beweis.}
    \[
        \glssymbol{wfk}_{\zvx} (e) = \pex(\kl{e}) = \txc{\pe(\zvx \in \kl{e}) = } \pe(\zvx=e) = \sum_{\substack{\oio \\ \zvx(\glssymbol{elementarereignis}) = e}} \glssymbol{wfk}(e) \geq 0
    \]
    z.z.\ Normierung
    \[
        \sum_{e \in E} \glssymbol{wfk}_{\zvx} (e) = \sum_{e \in E} \sum_{\substack{\oio :\\ \zvx(\glssymbol{elementarereignis}) = e}} \glssymbol{wfk} (\glssymbol{elementarereignis}) =  \sum_{\substack{\oio \\ e \in E :\\ \zvx(\glssymbol{elementarereignis}) = e}} \glssymbol{wfk}  (\glssymbol{elementarereignis}) = \sum_{\substack{\oio :\\ \zvx(\glssymbol{elementarereignis}) = e}} \glssymbol{wfk} (\glssymbol{elementarereignis}) = \sum_{\oio} \glssymbol{wfk} (\glssymbol{elementarereignis}) = 1
    \]
    \txc{(Ausführlicher auf Seite 31 im Skript 26)}
\end{lem}

\begin{kor}
    \textbf{Induzierter \gls{wraum}} \txc{(Skript 26 - Korollar 3.10.)}  \\
    $(E, \pex)$ ist ein \gls{wraum}.
\end{kor}

\begin{bsp}
    \txc{(Skript 26 - Beispiel 3.11.)}  \\
    Im Beispiel \autoref{bsp:n-faches-Werfen-einer-fairen-Muenze} und der Fortsetzung
    \autoref{bsp:n-faches-Werfen-einer-fairen-Muenze-2} erhalten wir als neuen, induzierten \gls{wraum}, je nach Wahl des
    Wertebereichs
    \begin{itemize}
        \item[(a)] $( \{0, \dots, n\}, \pe_{\sn} )$, wobei das \gls{wmaß} $\pe_{\sn}$ gegeben ist durch
        die \gls{wfk} $\{0, \dots, n\} \ni k \mapsto \binom{n}{k} 2^{-n}$,
        \item[(b)] $( \nz_0, \pe_{\sn} )$, wobei das \gls{wmaß} $\pe_{\sn}$ gegeben ist durch
        die \gls{wfk}
        \[
            \nz_0 \ni k \mapsto \begin{cases}
                                    \binom{n}{k} 2^{-n}, & \forall k \in \{0, \dots, n\} \\ 0 & sonst
            \end{cases}
        \]
    \end{itemize}
    Im folgenden wollen wir die Notation etwas vereinfachen, damit wir für Ereignisse, die durch eine
    \gls{zv} ausgedrückt werden, nicht immer wieder mühsam $\{\ele \in \eraum : \zvx(\ele) \in \dots \}$ schreiben
    müssen.
\end{bsp}


\section{Notationen}\label{sec:notationen} \txc{(2025)}
\[
    \pe(\zvx \in B) = \pe(\kl{\oio: \zvx(\glssymbol{elementarereignis}) \in B})
\]
Falls $B$ ein Intervall ist:
\[
    \pe(\zvx \geq b), \quad \pe(a < \zvx \leq b) , \dots
\]
\begin{defi}
    \textbf{Wertebereich} \txc{(2025)}  \label{def:wert-bild-e} \\
    Geg.:
    \[
        \zvx: \glssymbol{eraum}  \rightarrow E \txt{\gls{zv} auf } \wraum
    \]
    Dann heißt $E$ der Wertebereich und $\zvx(\glssymbol{eraum} )\coloneqq  \kl{\zvx(\glssymbol{elementarereignis}): \oio}$ heißt Bildbereich.
\end{defi}

\begin{lem}[Existenz einer \gls{zv}n mit vorgegebener \gls{verteilung})]
    \txc{(Skript 26 - Lemma 3.13.)}  \\
    Sei $(\tilde\eraum, \tilde\pe)$ ein beliebiger \gls{wraum}.
    Dann existieren ein \gls{wraum} $\wraum$ und eine \gls{zv} $\zvx$ auf $\wraum$
    derart, dass die \gls{verteilung} von $\zvx$ das vorgegebene \gls{wmaß} $\tilde\pe$ ist.
    d.h.\ es gilt $\pex = \tilde\pe$.

    \paragraph{Beweis.}
    Der Beweis ist verwirrend einfach.
    Wir können nämlich einfach folgende Wahl treffen:
    \[
        \eraum = \tilde\eraum, \quad \pe = \tilde\pe, \quad \zvx(\ele) = \ele \faoio
    \]
    Dann gilt für jede messbare Menge $B \subseteq \eraum$:
    \[
        \pex(B) = \pe(\{ \oio: \zvx(\ele) \in B \}) = \pe(B) = \tilde\pe(B)
    \]
    Dabei haben wir verwendet, dass $\zvx$ als Abbildung die Identität auf $\eraum$ ist und daher
    \[
        \{ \oio: \zvx(\ele) \in B \} = B
    \]
    gilt.
    Da $B \subseteq \eraum$  beliebig war, haben wir gezeigt, dass $\pe = \tilde\pe$ gilt.
\end{lem}

\begin{bsp}
    \defirefman{wert-bild-e} \txc{(2025)}
    \txo{
        \[
            \sn : \glssymbol{eraum} _n \rightarrow \kl{0,\dots,n}
        \]
        \[
            \sn : \glssymbol{eraum} _n \rightarrow \mathbb{N}_0
        \]
    }
\end{bsp}

\begin{bem}
    \defirefman{wert-bild-e} \txc{(2025)}  \\
    $\zvx(\glssymbol{eraum} ) \subset E$ gilt immer.
\end{bem}

\komm{Denn $\zvx(\glssymbol{eraum} ) $ ist ein $\oio$ und jedes \glssymbol{elementarereignis} ist ein Element aus $E$.}

\begin{bem}
    \defirefman{wert-bild-e} \txc{(2025)}  \\
    Gegeben ist ein \gls{wmaß} auf $\tilde {\glssymbol{eraum}}$, so existiert eine \gls{zv}, deren \gls{verteilung} genau dieses \gls{wmaß} ist. \\
    \noindent Sei $(\tilde {\glssymbol{eraum}} , \tilde\pe)$ ein beliebiger \gls{wraum}. \\
    $\zvx: \glssymbol{eraum}  \rightarrow \tilde{\glssymbol{eraum}} $ definiert durch:
    \begin{align*}
    {\glssymbol{eraum}}
        & \coloneqq  \tilde{\glssymbol{eraum}}  \\
        \pe&\coloneqq  \tilde\pe\\
        \zvx(\glssymbol{elementarereignis}) & = \glssymbol{elementarereignis} \txt{Identität} \faoio
    \end{align*}
    \[
        \pagebreak(\zvx \in B) = \pe(\oub{\kl{\oio: \zvx(\glssymbol{elementarereignis}) \in B}}{=B}) = \pe(B) = \tilde \pe(B)
    \]
    $\Rightarrow \zvx$ hat \gls{verteilung} $\tilde\pe$
\end{bem}

\komm{Zu jedem \gls{wmaß} auf einem Messraum existiert eine \gls{zv}, die genau dieses Maß als \gls{verteilung} hat. \\
aka. Für jede vorgegebene Wahrscheinlichkeitsverteilung kann man ein Zufallsexperiment konstruieren, das diese \gls{verteilung} als Ergebnis hat.}

\begin{defi}[\gls{verteilungsfunktion_zv}]
    \txc{(Skript 26 - Definition 3.14.)}  \\
    Seien $\wraum$ ein \gls{wraum} und $\zvx : \eraum \rightarrow \rz$ eine $\rz$-wertige \gls{zv}.
    Dann heißt die Funktion $\fxx : \rz \rightarrow [0,1]$, die gegeben ist durch
    \[
        \fxx(x) = \pe(\zvx \leq x) = \pex((-\infty,x]) \quad \forall x \in \rz
    \]
    die \gls{verteilungsfunktion_zv}.
\end{defi}

\begin{bem}
    Achtung, Verwechslungsgefahr!
    Wahrscheinlichkeitsfunktion $\wfk_{\zvx}$ und \gls{verteilungsfunktion_zv} $\fxx$ sind ver-
    schiedene Konzepte!
\end{bem}

\begin{bsp}[Fortsetzung des Münzwurf-Beispiels]
    \txc{(Skript 26 - Beispiel 3.15.)}  \\
    Wir werfen nun unsere faire Münze fünfmal, das heißt, wir betrachten die \gls{zv} $\zvs_5$.
    Die
    zugehörige \gls{verteilungsfunktion_zv} ist gegeben durch
    \[
        \fx_{\zvs_5}(x)
        =
        \pe(\zvs_5 \le x)
        =
        \begin{cases}
            0,
            &
            \text{für } x < 0,
            \\[0.5em]
            \pe(\zvs_5 \le 0)
            =
            \wfk_{\zvs_5}(0)
            =
            2^{-5}
            =
            \frac{1}{32},
            &
            \text{für } 0 \le x < 1,
            \\[0.5em]
            \pe(\zvs_5 \le 1)
            =
            \wfk_{\zvs_5}(0)
            +
            \wfk_{\zvs_5}(1)
            =
            2^{-5}
            +
            \binom{5}{1}2^{-5}
            =
            \frac{6}{32},
            &
            \text{für } 1 \le x < 2,
            \\[0.5em]
            \pe(\zvs_5 \le 2)
            =
            \sum_{i=0}^{2}\wfk_{\zvs_5}(i)
            =
            \fx_{\zvs_5}(1)
            +
            \wfk_{\zvs_5}(2)
            =
            \frac{1}{2},
            &
            \text{für } 2 \le x < 3,
            \\[0.5em]
            \pe(\zvs_5 \le 3)
            =
            \sum_{i=0}^{3}\wfk_{\zvs_5}(i)
            =
            \fx_{\zvs_5}(2)
            +
            \wfk_{\zvs_5}(3)
            =
            \frac{26}{32},
            &
            \text{für } 3 \le x < 4,
            \\[0.5em]
            \pe(\zvs_5 \le 4)
            =
            \sum_{i=0}^{4}\wfk_{\zvs_5}(i)
            =
            \fx_{\zvs_5}(3)
            +
            \wfk_{\zvs_5}(4)
            =
            \frac{31}{32},
            &
            \text{für } 4 \le x < 5,
            \\[0.5em]
            1,
            &
            \text{für } x \ge 5.
        \end{cases}
    \]
\end{bsp}

\begin{lem}[Eigenschaften von \gls{verteilungsfunktion_zv}en I: allgemeine Eigenschaften]
    \txc{(Skript 26 - Lemma 3.17.)}  \\
    Sei $\fxx$ eine \gls{verteilungsfunktion_zv}.
    Dann gelten:
    \begin{itemize}
        \item $\fxx$ ist monoton wachsend
        \item $\fxx$ ist rechtsstetig, d.h.: Für jedes $x_0 \in \rz$ und jede Folge $(x_n)_{n \in \nz}$ mit $x_n \searrow x_0$ gilt $\limnin \fxx (x_n) = \fxx (x_0)$
        \item $\lim_{x \rightarrow - \infty} \fxx =(x) = 0$ und $\limnin \fxx (x) = 1$
    \end{itemize}
\end{lem}

\begin{lem}[Rechnen mit \gls{verteilungsfunktion_zv}en I: allgemeine Rechenregeln]
    \txc{(Skript 26 - Lemma 3.17.)}  \\
    Sei $- \infty < a \leq b < \infty$.
    Für die \gls{verteilungsfunktion_zv} einer \gls{zv}n $\zvx$ gelten neben der definierenden Gleichheit (siehe (b)) auch
    die folgenden Rechenregeln:
    \begin{itemize}
        \item \( \pe (a < \zvx \leq b) = \fxx (b) - \fxx (a)\)
        \item \(\pe(\zvx \leq b) = \fxx (b)\)
        \item \(\pe(\zvx > a) = 1 - \fxx(a)\)
    \end{itemize}
\end{lem}


\section{\gls{zv}n mit abzählbarem Bildbereich}\label{sec:zufallsvariablen-mit-abzahlbarem-bildbereich}

%\paragraph{Hier fehlt was} % TODO

\begin{bsp}[Beispiel Geometrische \gls{verteilung}]
    \txc{(Skript 26 - Beispiel 3.18.)}  \\
    Wir werfen eine faire Münze bis das erste Mal Kopf fällt.
    In einem ersten Schritt begrenzen wir die Anzahl der Versuche und werfen maximal n-mal.
    Das offensichtliche Modell für das n-malige Werfen einer fairen Münze ist.
    \[
        \eraum_n = \{0,1\}^n \quad und \quad \wfk_n(\ele) = 2^{-n} \faoio
    \]
    Dabei bedeute $\ele_i = 1$, dass im i-ten Wurf Kopf fällt.

    Die einelementige Menge $A_n = \{( 0, 0,..., 0,1 )\} \subseteq \eraum_n$ beschreibt das Ereignis, dass im n-ten Wurf
    das erste Mal Kopf fällt.
    Es gilt $\pe_n(A_n) = \wfk_n(( 0, 0, \dots, 0,1 )) = 2^{-n}$.

    Wenn wir wirklich den Münzwurf wiederholen wollen, bis zum ersten Mal Kopf fällt, wissen wir nicht,
    welches n wir wählen sollten, da bei n Würfen stets eine positive Restwahrscheinlichkeit bestehen
    bleibt, dass wir keinen solchen Erfolg erzielen.

    Wir wählen daher einen anderen Ansatz, der die obigen Erkenntnisse nutzt.
    Sei $\eraum = \{1,2,\dots\} = \nz$.
    Dabei bedeute
    $\ele = k \in \eraum$, dass der erste Erfolg zur Zeit k eintritt.
    Gemäß der vorangegangenen
    Betrachtung wählen wir $\wfk(k) = 2^{-k}$ für alle $k \in \eraum = \nz$.

    Als nächstes überzeugen wir uns davon, dass $\wfk$ eine Wahrscheinlichkeitsfunktion ist.
    Offensichtlich ist $\wfk(k) > 0$ für alle $k \in \nz$, und es gilt
    \[
        \sum_{k=1}^{\infty} \wfk(k) =  \sum_{k=1}^{\infty} 2^{-k} = \frac{1}{2}  \sum_{l=0}^{\infty} 2^{-l}
    \]
    wobei wir den Summationsindex um 1 verschoben haben, indem wir $l = k - 1$setzen.
    Berechnen der
    geometrischen Reihe zeigt, dass
    \[
        \sum_{k=1}^{\infty} \wfk(k) = \frac{1}{2} \frac{1}{1- \frac{1}{2}} = 1
    \]
    Damit haben wir gezeigt, dass $\wfk$ eine Wahrscheinlichkeitsfunktion auf $\eraum = \nz$ ist. %, siehe Abbildung3.2.
    Eine N-wertige \gls{zv} $\zvx$, deren \gls{verteilung} durch $\wfk_{\zvx} = \wfk$ gegeben ist, nennen wir
    geometrisch verteilt.

    Der Bildbereich von $\zvx$ ist also $\nz$.
    Wir berechnen die \gls{verteilung}funktion $\fxx$ folgendermaßen:
    \begin{align*}
        \pe(\zvx \le x)
        =
        0,
        \text{für } x < 1,
        \fxx(k)
        =
        \pe(\zvx \le k)
        =
        \sum_{i=1}^{k}\wfk_{\zvx}(i)
        =
        \sum_{i=1}^{k}2^{-i}
        =
        1-\frac{1}{2^k},
        &
        \text{für } k\in\nz,
        \\
        \fxx(x)
        =
        \pe(\zvx \le x)
        =
        \pe(\zvx \le \lfloor x\rfloor)
        =
        \fxx(\lfloor x\rfloor)
        =
        1-\frac{1}{2^{\lfloor x\rfloor}},
        &
        \text{für } x\in[1,\infty)\setminus\nz.
    \end{align*}
\end{bsp}

% TODO $\wfk_{\zvx}$ definieren

\begin{bem}[Sprunghöhen der \gls{verteilungsfunktion_zv}]
    \txc{(Skript 26 - Bemerkung 3.19.)}  \\
    Am Beispiel der geometrischen \gls{verteilung} sehen wir in Abbildung 3.2 \txc{(siehe Skript 2026)}, dass die Sprunghöhen
    der \gls{verteilungsfunktion_zv} den Werten der Wahrscheinlichkeitsfunktion $\wfk_{\zvx}$ entsprechen.
    Das ist kein Zufall.
    Nimmt
    eine \gls{zv} $\zvx$ nur abzählbar viele Werte $x_1, x_2,\dots$ an, so gilt stets
    \begin{align*}
        \wfk_{\zvx}(x_i) ) & = \pe(\zvx = x_i) \\
        & = \pe(\zvx \leq x_i) - \pe(\zvx \le x_i) \\
        & = \pe(\zvx \leq x_i) - \pe(\zvx \leq x_{i-1}) \\
        & = \fxx (x_i) - \fxx(x_{i-1}) & \forall i \geq 2
    \end{align*}
    sowie
    \begin{align*}
        \wfk_{\zvx}(x_1) ) & = \pe(\zvx = x_1) \\
        & = \pe(\zvx \leq x_1) - \pe(\zvx \le x_1) \\
        & = \fxx(x_1) - 0 \\
        & = \fxx(x_1)
    \end{align*}
\end{bem}

Wir ergänzen nun unser Lemma über die Eigenschaften von \gls{verteilungsfunktion_zv}en um eine Aussage
über diskrete \gls{zv}n.

\begin{lem}[Eigenschaften von \gls{verteilungsfunktion_zv}en II: diskrete \gls{zv}]
    \txc{(Skript 26 - Lemma 3.20.)}  \\
    Sei $\fxx$ eine \gls{verteilungsfunktion_zv}.
    Dann gelten:
    \begin{itemize}
        \item $\fxx$ ist monoton wachsend
        \item $\fxx$ ist rechtsstetig, d.h.: Für jedes $x_0 \in \rz$ und jede Folge $(x_n)_{n \in \nz}$ mit $x_n \searrow x_0$ gilt $\limnin \fxx (x_n) = \fxx (x_0)$
        \item $\lim_{x \rightarrow - \infty} \fxx =(x) = 0$ und $\limnin \fxx (x) = 1$
        \item Nimmt die \gls{zv} $\zvx$ nur abzählbare viele Werte an, so ist $\fxx$ stückweise konstant.
    \end{itemize}
\end{lem}


\section{\gls{zv}n mit Dichte}\label{sec:zufallsvariablen-mit-dichte}

\begin{defi}[\dichte einer \gls{zv}]
    \txc{(Skript 26 - Definition 3.21.)}  \\
    Seien $(\eraum, \pe)$ ein Wahrscheinlichkeitsraum und $\zvx \colon \eraum \to \mathbb{R}$ eine
    $\mathbb{R}$-wertige \gls{zv}.
    Wir sagen, dass die \gls{zv} $\zvx$
    \emph{absolutstetig mit Dichte $f$} ist, wenn eine Dichte $f \colon \mathbb{R} \to [0,\infty)$
    existiert derart, dass
    \begin{equation}
        \label{eq:dichte-vf}
        \pex\bigl((-\infty, x]\bigr) = \pe(\zvx \leq x) = \int_{-\infty}^{x} f(u) du
        \qquad \forall\, x \in \mathbb{R}.
    \end{equation}
\end{defi}

\begin{bem}
    \txc{(Skript 26 - Bemerkung 3.22.)}  \\
    Umgekehrt können wir natürlich~\eqref{eq:dichte-vf} verwenden, um zu einer gegebenen Dichte $f$
    eine \gls{zv} $\zvx$ zu konstruieren, deren \gls{verteilung} $\pex$ die Dichte $f$ hat.

    Ohne Beweis merken wir an dieser Stelle an, dass~\eqref{eq:dichte-vf} das
    Wahrscheinlichkeitsmaß $\pex$ eindeutig festlegt.
\end{bem}

\begin{bsp}[Uniform verteilte \gls{zv}n]
    \txc{(Skript 26 - Beispiel 3.23.)}  \\
    Sei $f$ die Dichte der uniformen \gls{verteilung} auf einem Intervall $[a,b]$, das heißt, es gelte
    \[
        f(x) =
        \begin{cases}
            \frac{1}{b-a} & \text{für } x \in [a,b] ,\\
            0             & \text{sonst.}
        \end{cases}
    \]
    Sei $\zvx$ eine \gls{zv} mit Dichte $f$.
    Wir sagen, dass $\zvx$ \emph{uniform verteilt} ist.
    Wir wollen nun für diese \gls{zv} die \gls{verteilungsfunktion_zv} $\fxx$ berechnen.
    \begin{itemize}
        \item Als erstes betrachten wir $x < a$.
        Für alle $u < a$ gilt $f(u) = 0$, woraus
        \[
            \fxx(x) = \int_{-\infty}^{x} f(u) du = \int_{-\infty}^{x} 0 du = 0
        \]
        folgt.
        \item Sei nun $a \leq x \leq b$.
        Wir teilen das $\fxx(x)$ definierende Integral auf
        in ein Integral bis $a$ und ein weiteres Integral von $a$ bis $x$:
        \[
            \fxx(x) = \int_{-\infty}^{a} f(u) du + \int_{a}^{x} f(u) du
            = \fxx(a) + \int_{a}^{x} \frac{1}{b-a} du.
        \]
        Wir haben den ersten Summanden $\fxx(a) = 0$ bereits berechnet.
        Daher folgt
        \[
            \fxx(x) = \int_{a}^{x} \frac{1}{b-a} du = \frac{x-a}{b-a}.
        \]
        \item Es bleibt der Fall $x > b$ zu betrachten.
        Einerseits können wir wieder das Integral
        aufspalten, diesmal bei $b$, und erhalten
        \[
            \fxx(x) = \int_{-\infty}^{b} f(u) du + \int_{b}^{x} f(u) du = \fxx(b) + 0 = 1.
        \]
        Andererseits können wir auch schlau argumentieren: $\fxx$ ist monoton wachsend.
        Daraus folgt,
        dass $\fxx(x) \geqslant \fxx(b) = 1$ gelten muss.
        Da eine \gls{verteilungsfunktion_zv} durch $1$
        beschränkt ist, folgt bereits $\fxx(x) = 1$ für alle $x \geqslant b$.
    \end{itemize}
    Somit gilt
    \[
        \fxx(x) =
        \begin{cases}
            0                 & \text{für } x < a ,\\
            \dfrac{x-a}{b-a}  & \text{für } a \leq x \leq b ,\\
            1                 & \text{für } x > b.
        \end{cases}
    \]
\end{bsp}

\begin{bsp}[Exponentialverteilte \gls{zv}n]
    \txc{(Skript 26 - Beispiel 3.24.)}  \\
    Sei $f$ die Dichte der Exponentialverteilung mit Parameter $\alpha > 0$, das heißt, es gelte
    \[
        f(x) =
        \begin{cases}
            \alpha \eu^{-\alpha x} & \text{für } x > 0 ,\\
            0                      & \text{sonst.}
        \end{cases}
    \]
    Sei $\zvx$ eine \gls{zv} mit Dichte $f$.
    Wir sagen, dass $\zvx$ \emph{exponentialverteilt mit
    Parameter $\alpha$} ist.
    Wir wollen nun die \gls{verteilungsfunktion_zv} $\fxx$ von $\zvx$ bestimmen.
    \begin{itemize}
        \item Als erstes betrachten wir $x \leq 0$.
        Für alle $u \leq 0$ gilt $f(u) = 0$,
        woraus wieder
        \[
            \fxx(x) = \int_{-\infty}^{x} f(u) du = 0
        \]
        folgt.
        \item Sei nun $x > 0$.
        Wir teilen wieder das $\fxx(x)$ definierende Integral auf, diesmal in
        ein Integral bis $0$ und ein weiteres Integral von $0$ bis $x$:
        \[
            \fxx(x) = \fxx(0) + \int_{0}^{x} \alpha \eu^{-\alpha u} du
            = \int_{0}^{x} \alpha \eu^{-\alpha u} du.
        \]
        Der Integrand $h(u) = \alpha \eu^{-\alpha u}$ besitzt die Stammfunktion
        $H(u) = -\eu^{-\alpha u}$.
        Daher folgt
        \[
            \fxx(x) = H(u)\big|_{u=0}^{x} = 1 - \eu^{-\alpha x}.
        \]
    \end{itemize}
    Somit gilt
    \[
        \fxx(x) =
        \begin{cases}
            0                      & \text{für } x \leq 0 ,\\
            1 - \eu^{-\alpha x}    & \text{für } x > 0.
        \end{cases}
    \]
\end{bsp}

\begin{lem}[Rechnen mit \gls{verteilungsfunktion_zv}en II: \gls{zv} mit Dichte]
    \txc{(Skript 26 - Lemma 3.25.)}  \\
    Sei $-\infty < a \leq b < \infty$.
    Für die \gls{verteilungsfunktion_zv} einer \gls{zv}n $\zvx$
    mit Dichte $f$ gelten neben der definierenden Gleichheit (siehe (b)) auch die folgenden
    Rechenregeln.
    \begin{enumerate}[label=(\alph*)]
        \item $\pe(a < \zvx \leq b) = \displaystyle\int_{a}^{b} f(x) dx = \fxx(b) - \fxx(a)$,
        \item $\pe(\zvx \leq b) = \displaystyle\int_{-\infty}^{b} f(x) dx = \fxx(b)$,
        \item $\pe(\zvx > a) = \displaystyle\int_{a}^{\infty} f(x) dx = 1 - \fxx(a)$,
        \item \begin{align*}
                  \pe(a \leq \zvx < b) = \pe(a < \zvx \leq b) = \pe(a \leq \zvx \leq b)
                  &= \pe(a < \zvx < b) = \fxx(b) - \fxx(a) ,\\
                  \pe(\zvx \leq b) &= \pe(\zvx < b) = \fxx(b) ,\\
                  \pe(\zvx > a) &= \pe(\zvx \geqslant a) = 1 - \fxx(a).
        \end{align*}
    \end{enumerate}
\end{lem}

\begin{lem}[Eigenschaften von \gls{verteilungsfunktion_zv}en: finale Version]
    \txc{(Skript 26 - Lemma 3.26.)}  \\
    Sei $\fxx$ eine \gls{verteilungsfunktion_zv}.
    Dann gelten:
    \begin{enumerate}[label=(\alph*)]
        \item $\fxx$ ist monoton wachsend.
        \item $\fxx$ ist rechtsstetig, d.h.: Für jedes $x_0 \in \mathbb{R}$ und jede Folge
        $(x_n)_{n \in \mathbb{N}}$ mit $x_n \searrow x_0$ gilt
        $\lim_{n \to \infty} \fxx(x_n) = \fxx(x_0)$.
        \item $\lim_{x \to -\infty} \fxx(x) = 0$ und $\lim_{x \to \infty} \fxx(x) = 1$.
        \item Nimmt die \gls{zv} $\zvx$ nur abzählbar viele Werte an, so ist $\fxx$ stückweise
        konstant.
        \item Hat $\zvx$ eine Dichte, so ist $\fxx$ eine stetige Funktion.
        \item Hat $\zvx$ eine Dichte $f$ und ist $f$ stetig in einem Punkt $a \in \mathbb{R}$, so gilt
        $f(a) = \fxx'(a)$.
    \end{enumerate}
\end{lem}


\section{Mehrdimensionale \gls{zv}n und ihre gemeinsame \gls{verteilung}: Abzählbare Bildbereiche}
\label{sec:mehrdim-zv-abzaehlbar}

Für das folgende Beispiel ist die sogenannte charakteristische Funktion, die Ihnen bereits im Kontext
der \emph{Analysis} begegnet ist, nützlich.
Im Rahmen dieser Vorlesung werden wir stattdessen von
einer \emph{Indikatorfunktion} sprechen, weil der Begriff \emph{charakteristische Funktion} im
Bereich der Wahrscheinlichkeitstheorie eine andere Bedeutung hat.

\begin{defi}
    \textbf{Indikatorfunktion bzw. charakteristische Funktion} \txc{(Skript 26 - Defintion 3.27.)}  \\
    Sei $\eraum \neq \emptyset$ eine beliebige Menge.
    Für eine Teilmenge $B \subseteq \eraum$ heißt
    die Abbildung $\chaf_B \colon \eraum \to \{0,1\}$, die gegeben ist durch
    \[
        \chaf_B(\ele) =
        \begin{cases}
            1 & \text{für } \ele \in B ,\\
            0 & \text{für } \ele \notin B ,
        \end{cases}
    \]
    die \emph{Indikatorfunktion} bzw.
    die \emph{charakteristische Funktion der Menge $B$}.
\end{defi}

\begin{bsp}[Robin und Uli würfeln \dots]
    \txc{(Skript 26 - Beispiel 3.28.)}  \\
    Robin und Uli würfeln dreimal mit einem fairen Würfel.
    Robin gewinnt jedes Mal einen Euro, wenn
    eine Eins fällt, während Uli einen Euro gewinnt, wenn eine Sechs fällt.

    Wir modellieren das dreimalige Würfeln wie üblich mit $\eraum = \{1, \dots, 6\}^3$ und der
    Gleichverteilung $\pe$ auf $\eraum$.
    Die folgenden beiden \gls{zv}n beschreiben den
    Gewinn von Robin und jenen von Uli:
    \begin{align*}
        \zvx(\ele) &= \sum_{i=1}^{3} \chaf_{\{1\}}(\ele_i) \quad \forall\, \ele \in \eraum
        && \text{(Robins Gewinn)}\\
        \zvy(\ele) &= \sum_{i=1}^{3} \chaf_{\{6\}}(\ele_i) \quad \forall\, \ele \in \eraum
        && \text{(Ulis Gewinn)}
    \end{align*}
    Wir wollen beides gleichzeitig beschreiben.
    Daher sei $\zvy = (\zvx,\zvy)$.
    Die neue \gls{zv} $\zvy$
    hat den Wertebereich $\{0,1,2,3\}^2$.\footnote{Offensichtlich gilt in diesem Fall, dass der
    Bildbereich $\zvy(\eraum)$ eine echte Teilmenge von $\{0,1,2,3\}^2$ ist, weil es u.a.\ nicht
    möglich ist, dass sowohl Robin als auch Uli zweimal gewinnen.} Da der Wertebereich von $\zvy$
    zweidimensional ist, bietet es sich an, die Wahrscheinlichkeitsfunktion $\wfk_Z$ der \gls{verteilung}
    von $\zvy$ in Tabellenform anzugeben.

    Die Zeilensummen ergeben die \gls{verteilung} von $\zvx$.
    So gilt zum Beispiel, dass
    \[
        \pe(\zvx = 2) = \sum_{j=1}^{3} \pe(\zvx = 2, \zvy = j) = \sum_{j=1}^{3} \wfk_Z\bigl((2,j)\bigr)
        = 3 \Bigl(\tfrac16\Bigr)^2 \tfrac46 + 3 \Bigl(\tfrac16\Bigr)^3 + 0 + 0
        = \frac{15}{6^3} = 3 \Bigl(\tfrac16\Bigr)^2 \tfrac56.
    \]
    Rechts erkennen wir die Wahrscheinlichkeit, dass Robin in drei Runden zweimal gewinnt, wenn Uli
    gar nicht mitspielt.
    Ebenso ergeben die Spaltensummen die \gls{verteilung} von $\zvy$, und es gilt zum
    Beispiel, dass
    \[
        \pe(\zvy = 1) = \sum_{i=1}^{3} \pe(\zvx = i, \zvy = 1) = \sum_{i=1}^{3} \wfk_Z\bigl((i,1)\bigr)
        = \dots = \frac{75}{6^3}.
    \]
    Es bietet sich an, die Tabelle um die \gls{verteilung}en von $\zvx$ und $\zvy$ zu ergänzen:
    \[
        \renewcommand{\arraystretch}{1.6}
        \begin{array}{c|cccc|c}
            \wfk_Z((x,y)) & y = 0 & y = 1 & y = 2 & y = 3 & \wfk_\zvx(x) \\ \hline
            x = 0 & \bigl(\tfrac46\bigr)^3 & 3\,\tfrac16\bigl(\tfrac46\bigr)^2
            & 3\,\bigl(\tfrac16\bigr)^2\tfrac46 & \bigl(\tfrac16\bigr)^3
            & \bigl(\tfrac56\bigr)^3 \\
            x = 1 & 3\,\tfrac16\bigl(\tfrac46\bigr)^2 & 3!\,\bigl(\tfrac16\bigr)^2\tfrac46
            & 3\,\bigl(\tfrac16\bigr)^3 & 0
            & 3\,\tfrac16\bigl(\tfrac56\bigr)^2 \\
            x = 2 & 3\,\bigl(\tfrac16\bigr)^2\tfrac46 & 3\,\bigl(\tfrac16\bigr)^3 & 0 & 0
            & 3\,\bigl(\tfrac16\bigr)^2\tfrac56 \\
            x = 3 & \bigl(\tfrac16\bigr)^3 & 0 & 0 & 0 & \bigl(\tfrac16\bigr)^3 \\ \hline
            \wfk_\zvy(y) & \bigl(\tfrac56\bigr)^3 & 3\,\tfrac16\bigl(\tfrac56\bigr)^2
            & 3\,\bigl(\tfrac16\bigr)^2\tfrac56 & \bigl(\tfrac16\bigr)^3 & 1
        \end{array}
    \]
    Unten rechts sehen wir die Kontrolle: Die Summe über alle Felder im Inneren der Tabelle muss $1$
    betragen, genau wie die Summen der am Rand vermerkten Zeilen- bzw.\ Spaltensummen.
\end{bsp}

\begin{defi}
    \textbf{Gemeinsame \gls{verteilung} und Randverteilung von \gls{zv}n}
    \txc{(Skript 26 - Defintion 3.29.)}  \\
    Sei $(\eraum, \pe)$ ein Wahrscheinlichkeitsraum.
    \begin{enumerate}[label=(\alph*)]
        \item Seien $\zvx_1, \dots, \zvx_n$ \gls{zv}n, die auf $\eraum$ definiert sind.
        Dann heißt
        die \gls{verteilung} der $n$-dimensionalen \gls{zv}n $\zvx = (\zvx_1, \dots, \zvx_n)$ die
        \emph{gemeinsame \gls{verteilung} der \gls{zv}n $\zvx_1, \dots, \zvx_n$}.
        \item Ist andersherum eine $n$-dimensionale \gls{zv} $\zvx = (\zvx_1, \dots, \zvx_n)$ auf
        $\eraum$ gegeben, so heißt die \gls{verteilung} von $\zvx_i$ die \emph{$i$-te Randverteilung von $\zvx$}.
    \end{enumerate}
\end{defi}

\begin{bem}[Warnung!]
    \txc{(Skript 26 - Bemerkung 3.30.)}  \\
    Wir haben bereits in Beispiel 3.28 gesehen, wie wir aus einer gemeinsamen \gls{verteilung} von
    \gls{zv}n die Randverteilungen ausrechnen können.
    Leider gilt die Umkehrung nicht!
    Wir
    werden dies in einer der kommenden Vorlesungen genauer analysieren.

    Das Berechnen einer Randverteilung aus der gemeinsamen \gls{verteilung} von \gls{zv}n
    funktioniert im Fall von \gls{zv}n mit abzählbarem Bildbereich durch
    ``heraus summieren'' der restlichen Variablen.
\end{bem}

\begin{lem}[Bestimmen von Randverteilungen für \gls{zv}n mit abzählbarem Bildbereich]
    \txc{(Skript 26 - Lemma 3.31.)}  \\
    Seien $(\eraum, \pe)$ ein Wahrscheinlichkeitsraum und $\zvx_1, \dots, \zvx_n$ \gls{zv}n auf
    $(\eraum, \pe)$.
    Dabei sei jedes $\zvx_i$ eine $E_i$-wertige \gls{zv}.
    Ist $\wfk_X$ die
    Wahrscheinlichkeitsfunktion der gemeinsamen \gls{verteilung} der \gls{zv}n $\zvx_1, \dots, \zvx_n$,
    so ist für jedes $k \in \{1, \dots, n\}$ die Wahrscheinlichkeitsfunktion $\wfk_{\zvx_k}$ von $\zvx_k$
    gegeben durch
    \[
        \wfk_{\zvx_k}(x) = \sum_{u_1 \in E_1} \cdots \sum_{u_{k-1} \in E_{k-1}}
        \sum_{u_{k+1} \in E_{k+1}} \cdots \sum_{u_n \in E_n}
        \wfk_\zvx(u_1, \dots, u_{k-1}, x, u_{k+1}, \dots, u_n) \qquad \forall\, x \in E_k.
    \]
    \begin{proof}
        Um die Notation übersichtlich halten zu können, führen wir den Beweis für $n = 3$ und
        $k = 2$ und schreiben $\zvx = \zvx_1$, $\zvy = \zvx_2$, $\zvy = \zvx_3$ und entsprechend $x = u_1$, $y$ für
        die Werte von $\zvy = \zvx_2$ und $z = u_3$.
        Aus der Definition von $\wfk_X$ folgt, dass
        \[
            \wfk_\zvy(y) = \pe(\zvy = y) = \pe(\zvx \in E_1, \zvy = y, \zvy \in E_3)
            = \sum_{x \in E_1} \sum_{z \in E_3} \wfk_\zvx(x,y,z) \qquad \forall\, y \in E_2 ,
        \]
        also gilt $\wfk_\zvy(y) = \sum_{x \in E_1} \sum_{z \in E_3} \wfk_\zvx(x,y,z)$ für alle
        $y \in E_2$.
    \end{proof}
\end{lem}


\section{Mehrdimensionale \gls{zv}n und ihre gemeinsame \gls{verteilung}: Dichten}
\label{sec:mehrdim-zv-dichten}

\txc{VL vom 15. Mai 2026}

\begin{bsp}[Zufällige Punkte im Quadrat]
    \txc{(Skript 26 - Beispiel 3.32.)}  \\
    \emph{Wie groß ist die Wahrscheinlichkeit, dass ein im Einheitsquadrat $Q$ zufällig gewählter
    Punkt im Dreieck $D$ unter der Winkelhalbierenden liegt?}

    Das Quadrat $Q$ und das Dreieck $D$ sind gegeben durch
    \begin{align*}
        Q &= [0,1] \times [0,1] = \{(x,y) \in \mathbb{R}^2 \colon 0 \leq x \leq 1,\ 0 \leq y \leq 1\},\\
        D &= \{(x,y) \in Q \colon y \leq x\}.
    \end{align*}
    Intuitiv würden wir erwarten, dass $\pe(D) = \frac{|D|}{|Q|} = \frac12$ gilt, das heißt, wir
    erwarten, dass sich die Wahrscheinlichkeit als Verhältnis der Inhalte von Dreieck und Quadrat
    ergibt.

    Zum ``Aufwärmen'' beginnen wir damit, die Inhalte von Quadrat und Dreieck formal zu
    berechnen:
    \begin{align*}
        |Q| &= \int_{-\infty}^{\infty} \int_{-\infty}^{\infty} \chaf_Q(x,y) dx dy
        = \int_{0}^{1} \left( \int_{0}^{1} 1 dx \right) dy = \int_{0}^{1} 1 dy = 1 ,\\
        |D| &= \int_{-\infty}^{\infty} \int_{-\infty}^{\infty} \chaf_D(x,y) dx dy
        = \int_{0}^{1} \left( \int_{y}^{1} 1 dx \right) dy
        = \int_{0}^{1} (1-y) dy = 1 - \frac12 = \frac12.
    \end{align*}
    Dabei haben wir in der Berechnung des Inhalts des Dreiecks beim zweiten Gleichheitszeichen die
    Indikatorfunktion ersetzt durch das Einsetzen der Integralgrenzen, die die Indikatorfunktion
    vorschreibt.
    Um im Dreieck $D$ zu sein, müssen zwei Bedingungen erfüllt sein: Einerseits muss
    $x,y \in [0,1]$ gelten, andererseits muss auch $y \leq x$ gelten.
    Da das äußere Integral die
    Integrationsvariable $y$ hat, fordern wir dort nur, dass $y \in [0,1]$ gilt, was zum Integral
    $\int_0^1 \dots dy$ führt.
    (Die Grenzen des äußeren Integrals dürfen nicht von $x$ abhängen!)
    Das innere Integral hängt nun von $y$ ab.
    Wir behandeln $y \in [0,1]$ wie einen Parameter, wenn
    wir das innere Integral berechnen.
    Das innere Integral hat $x$ als Integrationsvariable, und $x$
    muss $x \in [0,1]$ sowie $y \leq x$ erfüllen.
    Wir fassen diese beiden Bedingungen zusammen
    und erhalten ein Integral der Form $\int_y^1 \dots dx$.
    Schließlich rechnen wir erst das
    innere Integral (als Funktion des Parameters $y$) und dann das äußere Integral aus.

    Alternativ können wir die Rechnung für den Inhalt des Dreiecks durch Vertauschen der
    Integrationsreihenfolge vereinfachen.
    Das Vertauschen ist möglich, weil die Indikatorfunktion
    eine messbare, nicht negative Funktion ist:
    \[
        |D| = \int_{-\infty}^{\infty} \int_{-\infty}^{\infty} \chaf_D(x,y) dx dy
        = \int_{-\infty}^{\infty} \int_{-\infty}^{\infty} \chaf_D(x,y) dy dx
        = \int_{0}^{1} \left( \int_{0}^{x} 1 dy \right) dx
        = \int_{0}^{1} x dx = \frac12.
    \]
    Da wir mit der uniformen \gls{verteilung} auf dem Quadrat arbeiten, gilt, wie erwartet,
    \[
        \pe(D) = \frac{\displaystyle\int_{-\infty}^{\infty} \int_{-\infty}^{\infty}
            \chaf_D(x,y) dx dy}
        {\displaystyle\int_{-\infty}^{\infty} \int_{-\infty}^{\infty}
            \chaf_Q(x,y) dx dy}
        = \frac12.
    \]
\end{bsp}

\begin{defi}
    \textbf{Mehrdimensionale Dichten}
    \txc{(Skript 26 - Defintion 3.33.)}  \\
    \begin{enumerate}[label=(\alph*)]
        \item Eine Funktion $f \colon \mathbb{R}^n \to [0,\infty)$ heißt \emph{$n$-dimensionale
        Dichte}, falls $f$ Riemann-integrierbar ist und die Normierungsbedingung
        \[
            \int_{-\infty}^{\infty} \cdots \int_{-\infty}^{\infty} f(x_1, \dots, x_n)
            dx_1 \dots dx_n = 1
        \]
        erfüllt ist.
        \item Eine $n$-dimensionale Dichte $f$ heißt \emph{gemeinsame Dichte der \gls{zv}n
            $\zvx_1, \dots, \zvx_n$}, wobei jedes $\zvx_i$ eine $E_i$-wertige \gls{zv} auf demselben
        Wahrscheinlichkeitsraum $(\eraum, \pe)$ ist, wenn
        \[
            \pe(\zvx_1 \leq a_1, \dots, \zvx_n \leq a_n)
            = \int_{-\infty}^{a_n} \cdots \int_{-\infty}^{a_1} f(x_1, \dots, x_n)
            dx_1 \dots dx_n
            \quad \forall\, a_1 \in E_1, \dots, a_n \in E_n
        \]
        gilt.
    \end{enumerate}
    Dabei verstehen wir unter $\{\zvx_1 \leq a_1, \dots, \zvx_n \leq a_n\}$ das Ereignis, dass
    alle $n$ Bedingungen simultan erfüllt sind, das heißt,
    \begin{align*}
        \{\zvx_1 \leq a_1, \dots, \zvx_n \leq a_n\}
        &= \{\ele \in \eraum \colon \zvx_1(\ele) \leq a_1 \text{ und } \dots \text{ und }
        \zvx_n(\ele) \leq a_n\}\\
        &= \{\zvx_1 \leq a_1\} \cap \dots \cap \{\zvx_n \leq a_n\}\\
        &= \bigcap_{i=1}^{n} \{\zvx_i \leq a_i\}.
    \end{align*}
\end{defi}

\begin{defi}
    \textbf{uniforme \gls{verteilung}}
    \txc{(Skript 26 - Defintion 3.34.)}  \\
    Sei $C \subset \mathbb{R}^n$ eine messbare Menge in $\mathbb{R}^n$ mit endlichem Inhalt, das
    heißt, es gelte $|C| = \lambda^n(C) < \infty$.
    Dann heißt die Funktion
    $f \colon \mathbb{R}^n \to [0,\infty)$, die durch
    \[
        f(x_1, \dots, x_n) = \frac{1}{|C|} \chaf_C\bigl((x_1, \dots, x_n)\bigr) =
        \begin{cases}
            \frac{1}{|C|} & \text{für } (x_1, \dots, x_n) \in C ,\\
            0             & \text{sonst,}
        \end{cases}
        \qquad \forall\, (x_1, \dots, x_n) \in \mathbb{R}^n
    \]
    definiert ist, \emph{Dichte der uniformen \gls{verteilung} auf $C$}.
\end{defi}

\begin{lem}[Berechnen von Randdichten]\label{lem:berechnen-randdichten}
    \txc{(Skript 26 - Lemma 3.35.)}  \\
    Sei $f \colon \mathbb{R}^n \to [0,\infty)$ eine $n$-dimensionale Dichte der gemeinsamen
    \gls{verteilung} der \gls{zv}n $\zvx_1, \dots, \zvx_n$.
    Für $k \in \{1, \dots, n\}$ ist die Funktion
    $f_k \colon \mathbb{R} \to [0,\infty)$, gegeben durch
    \[
        f_k(x) = \int_{-\infty}^{\infty} \cdots \int_{-\infty}^{\infty}
        f(u_1, \dots, u_{k-1}, x, u_{k+1}, \dots, u_n)
        dx_1 \dots dx_{k-1} dx_{k+1} \dots dx_n \qquad \forall\, x \in \mathbb{R} ,
    \]
    eine Dichte der \gls{zv}n $\zvx_k$.
    \begin{proof}
        Um von einer übersichtlicheren Notation profitieren zu können, führen wir den Beweis nur für
        $n = 2$ und $k = 1$.
        Wir schreiben $\zvx = \zvx_1$, $\zvy = \zvx_2$ sowie $x = x_1$ und $y = x_2$.
        Für beliebiges $a \in \mathbb{R}$ gilt
        \begin{align*}
            \pe(\zvx \leq a) &= \pe(\zvx \leq a, \zvy \in \mathbb{R})\\
            &= \int_{-\infty}^{\infty} \int_{-\infty}^{\infty}
            f(x,y) \chaf_{(-\infty,a]}(x) dx dy\\
            &= \int_{-\infty}^{\infty} \int_{-\infty}^{a} f(x,y) dx dy.
        \end{align*}
        Da $f$ eine nichtnegative, messbare Funktion ist, dürfen wir die Integrationsreihenfolge
        vertauschen.
        Damit folgt
        \[
            \pe(\zvx \leq a) = \int_{-\infty}^{a} \left( \int_{-\infty}^{\infty} f(x,y) dy
            \right) dx = \int_{-\infty}^{a} f_\zvx(x) dx
        \]
        mit $f_\zvx(x) = \int_{-\infty}^{\infty} f(x,y) dy$.
        Da $a \in \mathbb{R}$ beliebig war,
        folgt, dass $f_X$ eine Dichte von $\zvx$ ist.

        \emph{Es ist eine gute Übung für Sie, entlang der Schritte dieses Beweises den Beweis für
            $n = 3$ und $k = 2$ selbständig zu führen.}
    \end{proof}
\end{lem}

Im Vergleich zum Fall von Wahrscheinlichkeitsfunktionen erhalten wir hier die \gls{verteilung} von $\zvx_k$
aus der gemeinsamen \gls{verteilung} von $(\zvx_1, \dots, \zvx_n)$, indem wir die restlichen Variablen
``heraus integrieren'' anstelle von ``heraus summieren''.

\begin{defi}
    \textbf{Randdichte}
    \txc{(Skript 26 - Defintion 3.36.)}  \\
    Die Dichte $f_k$ von $\zvx_k$ gemäß Lemma \autoref{lem:berechnen-randdichten} nennen wir \emph{$k$-te Randdichte von
        $(\zvx_1, \dots, \zvx_n)$}.
\end{defi}

% TODO Beispielaufgaben

\newpage
